\section{What Delivery Buys}
\label{sec:buys}

If delivery does not decide accuracy once the facts are equal, what is the
governed contract for? Three things, each measured: it keeps cost flat as
the contract grows, it records which knowledge each answer used, and its
declared rules keep the agent from writing to the warehouse.

\subsection{Cost and scale}
\label{sec:scale}

\begin{table}[t]
\centering
\footnotesize
\setlength{\tabcolsep}{2.6pt}
\begin{tabular}{@{}llrrrrr@{}}
\toprule
& & 1x & 2x & 4x & 8x & all \\
\midrule
\multicolumn{2}{@{}l}{prompt, k tokens: governed} & 8 & 9 & 10 & 12 & 20 \\
\multicolumn{2}{@{}l}{\phantom{prompt, k tokens:} text} & 26 & 53 & 92 & 188 & 551 \\
\midrule
\mLuna{} & governed & 2.22 & 2.23 & 2.34 & 2.44 & 2.97 \\
& text & 1.75 & 2.93 & 4.05 & 7.57 & 19.10 \\
\mSol{} & governed & 42.0 & 42.3 & 44.3 & 46.4 & 57.4 \\
& text & 40.3 & 67.4 & 104.2 & --- & --- \\
\bottomrule
\end{tabular}
\caption{Experiment 6: dollars per run ($\times 10^{-3}$) as the contract
grows from one database (1x) to 22 (all), $k{=}3$ on 60 questions;
prompt sizes are \mLuna{}'s tokens.
Governed: the merged contract behind the governed tools. Text: the same
files in the prompt over plain tools. \mSol{}'s text arm at 8x and all was
dropped for cost by a recorded amendment. Accuracy is 98 to 100\% in every
cell.}
\label{tab:scale}
\end{table}

Experiment 4b found that looking knowledge up as needed saves input but
not money: one database's contract, about 25k tokens, fits in a cached
prompt, so loading it whole is nearly free. Experiment 6 grows the
contract until that should change. Experiment 4's three databases and 60
questions are kept, each with its base contract, and padded with real but
irrelevant contracts: the other two bases first, then 19 contracts the
author wrote for other LiveSQLBench databases, each namespaced under its
database's name, with the base at the middle. Five levels take whole
contracts, up to all 22 databases. The governed arm's prompt holds an
index that grows from 8k tokens to 20k; the text arm's holds the files,
from 26k to 551k (Table~\ref{tab:scale}). The common top level, the
largest that both arms fit on every model, is 4x: \mQwen{}'s window does
not fit the text arm beyond it.

\textbf{Size costs no accuracy, up to 4x.} The text arm at 4x is below
itself at 1x on 3 tasks and above on none; the 95\% interval of the mean
drop is $[+0.0, +1.3]$ points, inside the registered $\pm$5-point
equivalence bound. \mLuna{}'s text arm, with all 22 databases (551k
tokens) in its prompt, still scores 98.3\%, though that level was not
tested. With no drop to hold, the predictions that the governed arm's drop
is smaller, or that it drops at all, both fail. \mQwen{}'s governed arm
scores 100\% at 8x and at all 22 databases, where the text arm does not
fit its window. Within this design, then, the model is not lost in a long
prompt. The design's scope bounds that: the padding is other databases
with distinct names, and the questions sit at the ceiling at 1x. A
contract padded with neighbouring domains of one warehouse would be the
harder test.

\textbf{Cost is where the deliveries differ.} At one database the text arm
is as cheap or cheaper: its prompt is cached, and it takes 2.7 to 4.1
turns against the governed arm's 5. From 2x the governed arm is cheaper,
on 58 of 60 tasks at 4x on \mLuna{} and on all 60 on \mSol{}
($p{=}2{\times}10^{-15}$ and $9{\times}10^{-19}$): 1.7 and 2.4 times
cheaper there, and 6.4 times cheaper on \mLuna{} at 22 databases. The
governed arm's input doubles from 1x to all, as its index grows; the text
arm's grows 21-fold. Under one fixed price the result holds fully cached
and uncached, and uncached the governed arm is cheaper even at one
database. The traces show why the governed arm's cost stays flat: its
lookups never name a padding database or metric, and no query touches a
padding table. Padding reaches its context only through an unfiltered
metric listing, in 0 to 1.7\% of runs on the GPT models and up to 9\% on
\mQwen{}.

\subsection{A record of what each answer used}
\label{sec:audit}

Every governed lookup names a domain or metric and returns its
definition, so the trace records which knowledge each answer consulted
before its query was written. The text arm's trace shows the SQL, and
that the whole contract was in context, not which definition was relied
on. This analysis was not predicted. In experiment 6 the final SQL
contains a looked-up metric's expression verbatim in 91\% of \mLuna{}'s
governed runs, 98\% of \mSol{}'s and 87\% of \mQwen{}'s runs with a
lookup. A departure does not mark a wrong answer: in experiment 6 every
run that departed from its looked-up formula was right. What the record does is locate an
error's source. Of experiment 4b's 175 wrong governed answers, 174 are in
the database whose contract carries the author's misses, and 90 applied
the wrong definition verbatim; each names the contract entry that
produced it, to be fixed once for every later answer.
\textbf{The trust the governed contract offers is accountability, not
detection.}

\subsection{Governance}
\label{sec:governance}

The contract also declares forbidden operations, rendered into the prompt
and enforced by the tools' static checks. In the four single runs the
ungoverned arms submitted 166 mutating statements over 26 tasks and the
governed arms none (Fisher $p{=}3{\times}10^{-8}$ pooled). Over every
DABStep run since, the harness's integrity check on each run's working
copy found \textbf{50 corrupted copies in 29{,}866 ungoverned runs and
none in 22{,}946 governed ones}, on nine models; 44 of the 50 are
\mQwen{}'s. Two escalations in the single runs show the hazard: an agent
lost a temporary table to the tool's per-call connection, diagnosed it,
and switched to persistent tables (\appref{app:escalations}). The
enforcement layer never had to fire: the governed arms never attempted a
write, so the mechanism measured is deterrence by declared rule, and the
contract's text over the plain tools, whose prompt states the same rules,
also never corrupted a copy.
